% REVIEW DRAFT: not an input of the current published manuscript.
% This file is not proof certification; see geometry-audit-2026-11-21.json.

\section{Rank-one contraction and the second constant term}
\label{sec:sr-theta-two-step-contraction}

We prove the remaining constant-term implication by an adelic
rank-one argument. It uses continuity and polynomial growth;
no spectral expansion of a noncuspidal form is needed.

For a number field \(K\), write \(G=\mathrm{SL}_2(\mathbb A_K)\),
\(\Gamma=\mathrm{SL}_2(K)\), and
\(n(z)=\left(\begin{smallmatrix}1&z\\0&1\end{smallmatrix}\right)\).
Let \(u\) have component \(2\) at every infinite place and component
\(1\) at every finite place, and put
\(b=\operatorname{diag}(u,u^{-1})\).
The quotient \(K\backslash\mathbb A_K\) has probability Haar measure.

\begin{srlemma}[A contraction equation for an adelic constant term]
\label{lem:sr-theta-rank-one-contraction}
Let \(\phi\) be continuous and left \(\Gamma\)-invariant. Suppose that,
for an integer \(m\ge0\) and every compact \(\Omega\subset G\),
\[
 |\phi(b^k\kappa)|\le C_\Omega 2^{mk}
       \quad(k\ge0,\ \kappa\in\Omega).
 \tag{R1}\label{eq:sr-contraction-growth}
\]
Define \(c(g)=\int_{K\backslash\mathbb A_K}\phi(n(z)g)\,dz\).
If \(c\) is continuous and, for some \(|\lambda|<1\),
\[
 c(b^2g)-c(bg)=\lambda\bigl(c(bg)-c(g)\bigr)
                         \quad(g\in G),             \tag{R2}
 \label{eq:sr-contraction-recurrence}
\]
then \(c(bg)=c(g)\) for every \(g\).
\end{srlemma}

\begin{proof}
Put \(Dc(g)=c(bg)-c(g)\). Iteration gives
\[
 Dc(b^kg)=\lambda^kDc(g),\qquad
 c(b^kg)=c(g)+(1+\lambda+\cdots+\lambda^{k-1})Dc(g).
 \tag{R3}\label{eq:sr-contraction-geometric}
\]
Thus \(c(b^k\kappa)\) is bounded for \(k\ge0\) and
\(\kappa\) in a fixed compact set.

We specify the two elementary reduction facts used below.
Adelic reduction for \(\mathrm{SL}_2\) gives a compact \(\Omega_0\)
such that \(G=\Gamma\{b^k:k\ge0\}\Omega_0\).
Indeed a Siegel set consists of a compact unipotent factor,
a diagonal idele with norm bounded below, and a maximal compact
factor. The norm-one idele class group is compact: its finite part
is reduced to finitely many ideal classes, and Dirichlet's unit
lattice reduces its infinite logarithms to a compact parallelepiped.
Divide the remaining positive real norm into intervals
\([|u|^k,|u|^{k+1}]\). The bounded initial interval is absorbed
in \(\Omega_0\). This gives the asserted discrete-ray version of
the usual Siegel set.
\footnote{For the adelic reduction theorem and the definition of
these Siegel sets, see J.~Getz and H.~Hahn,
\emph{An Introduction to Automorphic Representations},
Theorem~2.7.2, author's text:
\url{https://sites.duke.edu/jgetz/files/2022/04/Graduate_Text.pdf}.}

The second fact is a useful quantitative form of additive
compactness. Given a sufficiently small neighborhood \(V\) of \(1\),
choose a nonnegative test function \(\eta\), smooth at infinity,
locally constant at finite places, with compact support in \(V\)
and integral one. For every compact \(\Omega\) there are a compact
\(\Omega'\) and a constant \(C\) such that, for every \(z\) and \(k\),
one can choose \(q\in K\) and \(v\in\Omega'\) with
\[
 n(-q)n(z)b^k=b^kv,\qquad
 |\eta(\kappa^{-1}v^{-1}y)-\eta(\kappa^{-1}y)|
           \le C2^{-k}\quad(\kappa\in\Omega,\ y\in G).
 \tag{R4}\label{eq:sr-contraction-bump}
\]
The union of the supports in \(y\) in this formula is compact,
independently of \(z,k,\kappa\).
Here are the details. The finite part of a compactly supported
locally constant function is invariant under some compact open
subgroup. Compactness of the finite projection of \(\Omega\)
allows us to choose an additive compact open subgroup \(U_f\)
so small that conjugating \(n(U_f)\) by any \(\kappa_f\) preserves
that invariance. Strong additive approximation chooses \(q\) with
\((z-q)_f\in U_f\); modulo the resulting fractional ideal its
infinite part can be placed in a fixed compact parallelepiped.
Consequently \(v=b^{-k}n(z-q)b^k\) has finite part in \(n(U_f)\)
and infinite off-diagonal entry bounded by \(C'2^{-2k}\).
The finite change of \(\eta\) is zero, and the mean-value estimate
for its infinite derivatives on a fixed compact set gives (R4),
even with \(2^{-2k}\) in place of \(2^{-k}\).
The same description proves compactness of all the supports.
A finite sum of products of infinite smooth functions and finite
locally constant functions gives these conclusions for \(\eta\).

Let \(\phi_\eta(g)=\int_G\phi(gy)\eta(y)\,dy\) and define \(c_\eta\)
similarly. Compact supports and compactness of the additive
quotient justify Fubini, so \(c_\eta\) is the unipotent average
of \(\phi_\eta\), and (R2) remains true after convolution.
Changing variables in the convolution, and using automorphy
and (R4), gives
\[
 \left|\phi_\eta(b^k\kappa)
          -\phi_\eta(n(z)b^k\kappa)\right|
       \le C_\Omega' 2^{(m-1)k}\qquad(m\ge1).
 \tag{R5}\label{eq:sr-contraction-growth-drop}
\]
To see all constants explicitly, write the difference as
\[
 \int_G \phi(b^ky)
       \{\eta(\kappa^{-1}y)
                 -\eta(\kappa^{-1}v^{-1}y)\}\,dy .
\]
Both terms are supported in one compact set. On that set (R1)
bounds the first factor by \(C_1 2^{mk}\), and (R4) bounds
the second by \(C_2 2^{-k}\). Its finite Haar volume supplies
the remaining constant. Averaging in \(z\), and using the
boundedness of \(c_\eta(b^k\kappa)\) from (R3), shows that
\(\phi_\eta\) has growth exponent \(m-1\).

We now induct on \(m\). For \(m=0\), the discrete-ray Siegel set
and automorphy make \(\phi\) globally bounded, say by \(M\).
Its unipotent average is bounded by \(M\), and (R3), applied
at \(b^{-k}g\), gives
\[
 |Dc(g)|\le 2M|\lambda|^k.
\]
Letting \(k\) tend to infinity proves \(Dc=0\).
For \(m\ge1\), every sufficiently small neighborhood admits
the bump just constructed. Equation (R5) and the induction
hypothesis give \(Dc_\eta=0\).
Choose such nonnegative bumps of integral one with supports
shrinking to \(1\). Continuity of \(c\) implies
\(c_\eta(g)\to c(g)\) and \(c_\eta(bg)\to c(bg)\).
Therefore \(Dc=0\) again. All convolutions are over compact sets;
no integrability over \(\Gamma\backslash G\) was assumed.
\end{proof}

We apply this lemma to the actual theta constant terms.
Write \(Y=p+q\), \(d=a+a'\), and assume
\[
 a'>2Y+1,\qquad \delta=(a-a')/2>2Y+1.
 \tag{R6}\label{eq:sr-two-step-range}
\]
Let \(\Sigma\) be the span of the large and small theta-period
images. Its line constant term has the two distinct Levi laws
\[
        \alpha(t)|t|_E^{a'/2},
        \qquad \beta(t)|t|_E^{a/2}.                \tag{R7}
 \label{eq:sr-two-step-upper-laws}
\]
We call the corresponding functions \(F_\alpha,F_\beta\).
They are unique: restrict the Levi to an idele of norm greater
than one, whose two eigenvalues have different absolute values.
The uniqueness also makes these components linear in \(F\).

\begin{srlemma}[The next \(\beta\)-component vanishes]
\label{lem:sr-theta-two-step-kill}
Suppose \(p\ge1\), \(F\in\Sigma\), and \(F_\beta=0\).
Remove a hyperbolic plane from the large first space and from
\(W\), obtaining the actual lower pair
\[
     (b,a';p-1,q),\qquad b=a-2.
\]
Assume the already proved proportionality and nonvanishing for
that lower pair and for its further rank reductions.
For every maximal-compact \(k\), realize
\(y\mapsto F_\alpha(\operatorname{lev}(y)k)\) as an element of
the lower common image \(\Sigma_1\), using
Proposition~\ref{prop:sr-theta-small-alpha-realization}.
Its line constant term on \(U(W_1)\) has zero \(\beta\)-component.
\end{srlemma}

\begin{proof}
If \(W_1\) has no isotropic line the assertion has no component
to test. Otherwise choose two mutually orthogonal hyperbolic
planes in \(W\), with isotropic vectors \(x_1,x_2\).
The first removal uses \(x_1\), and the next removal uses \(x_2\).
For every \(x\in U(W)(\mathbb A_F)\) the function
\[
           y\longmapsto
                 F_\alpha(\operatorname{lev}(y)x)
\]
admits a lower splitting into functions \(A_x,C_x\) with laws
\[
       \alpha(t)|t|_E^{a'/2},
       \qquad \beta(t)|t|_E^{b/2}.                \tag{R8}
 \label{eq:sr-two-step-lower-laws}
\]
The assertion for arbitrary \(x\), and not just compact \(x\),
is needed here. Write \(x=n\,m(t)\operatorname{lev}(g_1)k\)
by Iwasawa decomposition for the first line parabolic.
The first constant term is left \(n\)-invariant; \(\operatorname{lev}(y)\)
normalizes this unipotent group and commutes with \(m(t)\).
The asserted function is therefore the right translate by \(g_1\)
of the compact-\(k\) lower function, multiplied by
\(\alpha(t)|t|^{a'/2}\). This transports its lower splitting to
all \(x\). No polynomial test-space invariance under arbitrary
real group elements is being assumed.

Set \(c(x)=A_x(1)+C_x(1)\). Embed
\(\mathrm{SL}_2(\mathbb A_E)\) in \(U(W)(\mathbb A_F)\)
by its standard action on \(E x_1+E x_2\), inverse conjugate
transpose action on the dual isotropic plane, and identity on
the orthogonal complement. Write \(\iota\) for this embedding.
For the idele \(u\) in the preceding lemma, put \(r=|u|_E>1\)
and \(h=\iota(\operatorname{diag}(u,u^{-1}))\).
Equations (R7)--(R8) give
\[
 \begin{split}
 c(hx)&=A_x(1)+\lambda C_x(1),\\
 c(h^2x)&=A_x(1)+\lambda^2 C_x(1),\\
 \lambda&=\alpha(u)\beta(u^{-1})\,r^{(a'-b)/2},
                         \qquad |\lambda|<1.
 \end{split}                                      \tag{R9}
 \label{eq:sr-two-step-eigenvalue}
\]
For example the first character on the outer line contributes
\(\alpha(u)r^{a'/2}\). On the inner line the \(A\)-term contributes
its inverse, while the \(C\)-term contributes
\(\beta(u^{-1})r^{-b/2}\). This proves both formulas with the
actual, possibly nontrivial, unitary characters.
The strict inequality follows from \(b>a'\), a consequence of (R6).

For completeness we identify the average to which the rank-one
lemma applies. Let \(N_2\) be the radical of the parabolic of the
isotropic plane \(E x_1+E x_2\), and define, for fixed \(x\),
\[
 \phi_x(g)=\int_{N_2(F)\backslash N_2(\mathbb A_F)}
                         F(n\iota(g)x)\,dn .
\]
The quotient is compact and has probability measure.
The group \(\mathrm{SL}_2(E)\) normalizes \(N_2\), preserves
its rational subgroup, and preserves the probability measure.
Thus \(\phi_x\) is a continuous automorphic function for
\(\mathrm{SL}_2(E)\). It has the growth in (R1) because \(F\)
has polynomial adelic matrix height growth and the integration
uses a compact set of representatives.

Its upper-unipotent average is exactly \(g\mapsto c(\iota(g)x)\).
Here is the group calculation behind this assertion.
The unipotent group for the flag \(E x_1\subset E x_1+E x_2\)
has two decompositions:
\[
 N_{\rm flag}=N_{x_1}\rtimes N_{x_2,W_1}
             =N_2\rtimes\iota(N_{\mathrm{SL}_2}).
\]
In the block basis
\((x_1,x_2;\,W_2;\,y_1,y_2)\), the second factor on the right
is upper triangular unipotent on the first two coordinates
and inverse transpose on the last two; \(N_2\) has identity
on both of these diagonal two-by-two blocks.
This identifies both groups and their kernels directly.
All three rational unipotent quotients have probability Haar
measure, so integration in stages in either order is Fubini
on compact quotients. The first order gives the first line
constant term and then its lower line constant term,
namely \(A_x(1)+C_x(1)\); the second gives the asserted
rank-one average.

Equation (R9) supplies (R2) at every \(g\).
Lemma~\ref{lem:sr-theta-rank-one-contraction} consequently gives
\(c(hx)=c(x)\). Comparing this equality with (R9) yields
\((\lambda-1)C_x(1)=0\), hence \(C_x(1)=0\).
Replace \(x\) by \(\operatorname{lev}(y)k\) to obtain
the vanishing of the entire lower \(\beta\)-component, as claimed.
\end{proof}

\begin{srproposition}[One zero component forces the other one to vanish]
\label{prop:sr-theta-two-step-zero}
Assume proportionality and nonvanishing for every smaller
signature sum than \(Y\), with the actual rank reductions above.
If \(q\ge1\) and \(F\in\Sigma\) has \(F_\beta=0\), then
\(F_\alpha=0\).
For \(q=0\) it is enough to assume \(F_\alpha=0\):
in that case \(F_\beta=0\) follows from the boundary vanishing
already proved.
\end{srproposition}

\begin{proof}
If \(p=0\), the boundary vanishing
Corollary~\ref{cor:sr-theta-boundary-component-zero} already gives
\(F_\alpha=0\). Suppose \(p\ge1\). The compact-\(k\) restrictions
of \(F_\alpha\) lie in the lower common image \(\Sigma_1\).
The local constant-term descent proved in
Proposition~\ref{prop:sr-theta-constant-term-descent} shows that,
as the upper common-quotient vector varies in a submodule,
their span is a submodule of \(\Sigma_1\).
At finite places this is ordinary group equivariance.
At infinity it is equivariance for the lower Lie algebra
and maximal compact, which preserve the polynomial test spaces
in these adapted Cartan coordinates. These are sufficient for
the irreducibility argument; no real group translate is asserted
to preserve the polynomial test space.

Apply this observation to the module generated by \(F\).
The identity \(F_\beta=0\) is stable under right translation
(and under the Lie algebra and maximal compact on polynomial
vectors), since it is an identity of component functions on
the whole upper group. Thus all the lower restrictions in
this module have zero \(\beta\)-component by the preceding lemma.
By lower proportionality and Factor, \(\Sigma_1\) is one
irreducible common-quotient image, and it is nonzero.
If the lower restrictions had a nonzero span, this span
would be all of \(\Sigma_1\).

When \(q\ge1\), however, \(\Sigma_1\) contains a vector with
nonzero \(\beta\)-component. Its large-side zero evaluation
is onto the further lower polynomial test space. Choose
a test vector with nonzero period in the \(B\)-reduction
of \((b,a';p-1,q)\), which exists by the smaller-signature
hypothesis. Lift it by this zero evaluation. The exact
large-side \(\beta\) constant-term formula is a nonzero
geometric constant times that lower period.
This contradicts zero \(\beta\)-component on \(\Sigma_1\).
Consequently every compact-\(k\) lower restriction is zero.
Iwasawa decomposition and the first Levi law give
\(F_\alpha=0\) on the whole group.

For \(q=0\), the assertion needed for the induction is
instead the final assertion in the proposition:
the \(\beta\)-component vanishes on both original images
by Corollary~\ref{cor:sr-theta-boundary-component-zero}.
Once the \(\alpha\)-components have been matched it
therefore vanishes on their difference as well.
\end{proof}

In the case \(p,q\ge1\), the further \(B\)-reduction just used
has signature sum \(Y-2\), so this argument uses only earlier
induction cases. The line-to-cusp implication and the proved
cuspidal orthogonality then kill the difference whose chosen
component was matched. The remaining work is to fix the
adapted models before the finite packets and to identify
that chosen component's nonzero scalar.
